W dobie intensywnego rozwoju systemów informatycznych, aplikacji webowych oraz urządzeń Internetu Rzeczy (IoT), kluczowym wyzwaniem staje się wydajne przetwarzanie i przechowywanie ogromnych wolumenów danych w czasie rzeczywistym. Tradycyjne relacyjne bazy danych (RDBMS) oparte na języku SQL, choć gwarantują silną spójność i dojrzałość technologiczną, często okazują się niewystarczające pod kątem skalowalności poziomej i opóźnień w dostępie do danych. Odpowiedzią na te ograniczenia stał się rozwój systemów NoSQL (ang. \emph{Not Only SQL}), a w szczególności baz danych typu klucz-wartość (ang. \emph{key-value}). Charakteryzują się one prostym modelem danych, minimalnymi opóźnieniami oraz wysoką przepustowością, stanowiąc fundament systemów buforowania (caching), zarządzania sesjami oraz brokerów wiadomości.

Rynek systemów klucz-wartość jest bardzo zróżnicowany. Dostępne są zarówno rozwiązania pamięciowe (in-memory), bazy wbudowane (embedded), jak i zaawansowane systemy rozproszone. Różnią się one od siebie architekturą, modelami współbieżności, mechanizmami zapewniania trwałości danych oraz modelami licencjonowania. Zrozumienie tych różnic i ich praktycznego wpływu na wydajność w określonych środowiskach systemowych jest kluczowe dla inżynierów projektujących nowoczesne systemy informatyczne.

Niniejsza praca jest poświęcona szczegółowej analizie porównawczej ośmiu wyselekcjonowanych baz danych tego typu, reprezentujących odmienne podejścia architektoniczne: Redis, Memcached, LevelDB, RocksDB, Couchbase, Tarantool, FoundationDB oraz BerkeleyDB.

\section{Cel i zakres pracy}
Głównym celem pracy jest przeprowadzenie analizy porównawczej wydajności wybranych nierelacyjnych baz danych typu klucz-wartość w środowisku zwirtualizowanym WSL 2 oraz natywnym systemie Linux. 

Zakres pracy obejmuje:
\begin{itemize}
    \item Szczegółowe studia literaturowe dotyczące teorii baz danych NoSQL, twierdzenia CAP oraz architektury porównywanych systemów.
    \item Przygotowanie zunifikowanego środowiska testowego przy użyciu technologii kontenerowej Docker oraz Docker Compose.
    \item Opracowanie dedykowanej aplikacji benchmarkującej w języku Python 3.14, implementującej interfejsy dla wszystkich ośmiu baz danych.
    \item Przeprowadzenie serii eksperymentów badawczych mierzących przepustowość i opóźnienia przy operacjach CRUD, mixed workloads, operacjach na dokumentach JSON oraz strukturach listowych.
    \item Porównanie narzutu wirtualizacji (WSL 2 vs natywny Linux).
\end{itemize}

\section{Zakres badanych technologii}
Analizie poddano osiem systemów bazodanowych, które podzielono na trzy główne grupy architektoniczne:
\begin{itemize}
    \item \textbf{Systemy pamięciowe (In-Memory):} Redis (silnik struktur danych), Memcached (klasyczny bufor pamięciowy) oraz Tarantool (serwer aplikacji w pamięci RAM z integracją LuaJIT).
    \item \textbf{Systemy osadzone (Embedded Library):} LevelDB (stworzona przez Google biblioteka oparta na LSM-Tree), RocksDB (rozwinięty przez firmę Meta silnik zoptymalizowany pod dyski SSD) oraz BerkeleyDB (dojrzała biblioteka transakcyjna ACID od firmy Oracle).
    \item \textbf{Systemy rozproszone i wielomodelowe:} Couchbase Server (zorientowana na pamięć baza dokumentowa korzystająca z języka SQL++) oraz FoundationDB (rozproszona baza transakcyjna ACID rozwijana pod patronatem Apple).
\end{itemize}

\section{Hipotezy i pytania badawcze}
W pracy sformułowano następujące pytania badawcze:
\begin{enumerate}
    \item Jak model wdrożenia (klient-serwer vs baza osadzona) wpływa na skalowalność i opóźnienia przy rosnącej liczbie współbieżnych procesów klienckich?
    \item Jaki jest rzeczywisty narzut wydajnościowy związany z emulacją operacji na strukturach złożonych (JSON, tablice) w warstwie aplikacji w porównaniu z ich natywną obsługą po stronie silnika bazy danych?
    \item W jakim stopniu technologia wirtualizacji WSL 2 degraduje wydajność operacji wejścia/wyjścia na dysk (I/O) w porównaniu do natywnego systemu Linux?
\end{enumerate}

Na podstawie pytań badawczych postawiono następujące hipotezy:
\begin{itemize}
    \item \textbf{Hipoteza 1:} Systemy pamięciowe (Redis, Memcached, Tarantool) zapewnią najniższe opóźnienia i najwyższą przepustowość dla operacji odczytu, jednak ich skalowalność pionowa powyżej liczby fizycznych rdzeni procesora będzie ograniczona współdzieleniem zasobów w technologii SMT.
    \item \textbf{Hipoteza 2:} Bazy osadzone (RocksDB, BerkeleyDB) wykażą wyższą wydajność zapisu w środowiskach wieloprocesowych w porównaniu do jednowątkowych systemów klient-serwer z uwagi na brak narzutu protokołu sieciowego.
    \item \textbf{Hipoteza 3:} Wykonywanie modyfikacji dokumentów JSON przy użyciu emulacji w języku Python będzie wolniejsze o co najmniej rząd wielkości w porównaniu z natywnym przetwarzaniem w silnikach Redis czy Couchbase z powodu kosztów serializacji i przesyłu danych.
    \item \textbf{Hipoteza 4:} Narzut warstwy wirtualizacji w systemie WSL 2 spowoduje spadek wydajności baz zorientowanych dyskowo (LSM-Tree) podczas operacji zapisu z wymuszonym WAL w stosunku do natywnego systemu Linux.
\end{itemize}

\section{Metodyka badań}
Badania oparto na autorskim skrypcie benchmarkującym napisanym w języku Python w wersji 3.14. W celu wyeliminowania wpływu ograniczeń interpretera (blokada GIL), obciążenie generowane jest za pomocą procesów systemowych (od 1 do 16 procesów roboczych). Wszystkie bazy danych uruchamiane są w odizolowanych kontenerach Docker w systemach operacyjnych Ubuntu 24.04 LTS (zarówno natywnie, jak i w WSL 2). W trakcie testów kontrolowane są parametry takie jak: liczba przydzielonych rdzeni CPU dla bazy danych (1, 2 lub 4 rdzenie), rozmiar rekordu (128 B oraz 4 KB) oraz wolumen danych (10k, 100k, 1M). Rejestrowane metryki obejmują przepustowość, centyle opóźnień oraz zużycie zasobów sprzętowych hosta.

\section{Struktura pracy}
Struktura pracy została podzielona na następujące części:
\begin{itemize}
    \item \textbf{Rozdział 1 (Wstęp):} zawiera uzasadnienie wyboru tematu, sformułowanie pytań badawczych i hipotez oraz ogólny zarys metodyki.
    \item \textbf{Rozdział 2 (Przegląd teoretyczny):} przedstawia teorię systemów NoSQL, ich klasyfikację ze szczególnym uwzględnieniem baz typu klucz-wartość oraz zagadnienia spójności danych.
    \item \textbf{Rozdział 3 (Charakterystyka porównywanych systemów):} zawiera szczegółową analizę architektoniczną oraz analizę modeli licencjonowania ośmiu badanych systemów bazodanowych.
    \item \textbf{Rozdział 4 (Środowisko testowe i metodyka testów):} opisuje konfigurację sprzętowo-programową stanowiska, budowę kodu testowego oraz szczegółowe scenariusze badawcze.
    \item \textbf{Rozdział 5 (Wyniki eksperymentów):} zawiera prezentację wyników pomiarów w postaci tabel i wykresów dla środowisk Linux oraz WSL 2.
    \item \textbf{Rozdział 6 (Dyskusja wyników):} poświęcony jest analizie uzyskanych rezultatów pod kątem postawionych hipotez oraz sformułowaniu zaleceń inżynierskich.
    \item \textbf{Rozdział 7 (Podsumowanie):} podsumowuje zrealizowane badania, ocenia stopień osiągnięcia celów oraz wskazuje kierunki dalszych prac.
\end{itemize}
